\documentclass[14pt]{extarticle}
\usepackage{preamble}
\geometry{letterpaper, landscape, margin=0.5in}
\usepackage{qrcode}
\renewcommand{\answer}[1]{}
\setlength{\parskip}{4pt}

\begin{document}
\pagestyle{empty}

\daybanner{Unit 3 \textperiodcentered\ Topic 3.3 \textperiodcentered\ TODAY'S PROBLEM}{Six Activations in Twelve Users}

\noindent\begin{minipage}[t]{0.57\linewidth}
\vspace{0pt}
Suppose a developer's model says $p_0=0.20$ of new app users activate voice commands. An SRS of 12 users includes 6 activations, so $\hat p=0.50$. The chart shows 200 simulated samples of 12 independent users under this model.\par\smallskip \textbf{Mara:} ``A normal model for $\hat p$ gives $P(\hat p\ge 0.50)\approx0.005$, so that is the chance probability I would report.''\par \textbf{Jules:} ``The simulated distribution is not very normal. I would estimate the chance probability directly from its tail.''\par
\end{minipage}\hfill
\begin{minipage}[t]{0.40\linewidth}
\vspace{0pt}
\begin{center}
\begin{tikzpicture}
\begin{axis}[
  width=0.98\linewidth,height=1.75in,
  ybar,bar width=10pt,xmin=-0.6,xmax=6.6,ymin=0,ymax=70,
  xtick={0,1,2,3,4,5,6},ytick={0,20,40,60},
  xlabel={Activations in a sample of 12},ylabel={Simulated samples},
  tick label style={font=\small},label style={font=\small},
  nodes near coords,nodes near coords style={font=\scriptsize},
  axis lines*=left
]
\addplot[draw=dayblue,fill=dayblue!20] coordinates {(0,14) (1,40) (2,58) (3,47) (4,25) (5,12) (6,4)};
\end{axis}
\end{tikzpicture}
\end{center}
\end{minipage}

\begin{firsttakebox}{FIRST TAKE (5 min, on your sheet)}
Whose method would you trust more, Mara's smooth curve or Jules's repeated trials? Give one reason from the picture.
\end{firsttakebox}

\begin{description}[leftmargin=1.15in, labelwidth=1in, itemsep=2pt, topsep=2pt]
  \item[\dokbadge{1}~(a)] State the most common number of activations in the simulated samples and describe the distribution as roughly symmetric, skewed left, or skewed right.
  \item[\dokbadge{2}~(b)] Use the chart to estimate $P(\hat p\ge0.50)$ under $p_0=0.20$. Then calculate $np_0$ and $n(1-p_0)$ and assess the expected-count condition for a normal approximation.
  \item[\dokbadge{3}~(c)] In 3--4 sentences, choose a method using the simulated tail and either the shape or expected counts. Apply the $5\%$ evidence rule. Explain how reporting 0.005 instead of 0.020 would change a reader's impression.
\end{description}

\vfill
\noindent\begin{minipage}{0.84\linewidth}
\textbf{Video + follow-along:} \href{https://robjohncolson.github.io/apstats-live-worksheet/u6_lesson1-2_live.html}{\texttt{\detokenize{u6_lesson1-2_live.html}}} (scan the code)\par
\textbf{Finish (a)--(c) after the video. Turn in your sheet.}
\end{minipage}\hfill
\qrcode[height=0.75in]{https://robjohncolson.github.io/apstats-live-worksheet/u6_lesson1-2_live.html}

\end{document}
